\documentclass[12pt,a4paper]{article}
\usepackage[utf8]{inputenc}
\usepackage[T2A]{fontenc}
\usepackage[russian]{babel}
\usepackage{amsmath,amssymb}
\usepackage[left=25mm,right=25mm,top=20mm,bottom=20mm]{geometry}
\usepackage{enumitem}

\title{Функции псевдодавления и псевдовремени\\
для многофазной неизотермической фильтрации}
\author{}
\date{}

\begin{document}
\maketitle
\tableofcontents
\newpage

%=============================================================================
\section{Введение}\label{sec:intro}
%=============================================================================

Задача: выяснить, можно ли обобщить формализм псевдодавления
и псевдовремени на случай \emph{неизотермической} двухфазной
фильтрации (нефть~+~свободный газ) в рамках модели Black~Oil.

Классическая цепочка обобщений:
\begin{enumerate}[nosep]
  \item Al-Hussainy, Ramey, Crawford (1966): псевдодавление $m(p)$
        для однофазного реального газа --- линеаризация
        пространственного оператора;
  \item Agarwal (1979): псевдовремя $t_a(t)$ --- линеаризация
        временного оператора;
  \item Краснов (б/г): многофазное псевдодавление $m^k(p)$ и
        псевдовремя $t_{ao}(t)$ для изотермической системы Black~Oil;
  \item \textbf{Данный документ}: попытка распространить подход на
        систему с уравнением энергии.
\end{enumerate}

%=============================================================================
\section{Однофазный реальный газ (Al-Hussainy \& Agarwal)}\label{sec:single}
%=============================================================================

\subsection{Исходное уравнение}\label{ssec:single_eq}

Изотермическая фильтрация реального газа в пористой среде.
Уравнение неразрывности + закон Дарси:
\begin{equation}\label{eq:gas_cont}
  \frac{\partial}{\partial t}\bigl(\phi\rho_g\bigr)
  + \mathrm{div}\bigl(\rho_g\,\vec{u}_g\bigr) = 0,
  \qquad
  \vec{u}_g = -\frac{k}{\mu_g(p)}\nabla p.
\end{equation}
Плотность газа $\rho_g = pM/(Z(p)RT)$. Пористость~$\phi$ считаем
постоянной (сжимаемость породы мала по сравнению со сжимаемостью газа).
Поскольку рассматривается изотермическая фильтрация однокомпонентного
газа, молярная масса~$M$ и температура~$T$ --- константы. При подстановке
$\rho_g$ в уравнение неразрывности множитель $M/(RT)$ сокращается
в обеих частях, и уравнение принимает вид:
\begin{equation}\label{eq:gas_pde}
  \frac{\partial}{\partial t}\left(\frac{p}{Z}\right)
  = \mathrm{div}\left(\frac{kp}{\mu_g Z}\,\nabla p\right).
\end{equation}
Уравнение нелинейно: коэффициент $p/(\mu_g Z)$ зависит от давления.

\subsection{Замена переменной: псевдодавление}\label{ssec:pseudo_p}

Введём функцию псевдодавления (Al-Hussainy et al., 1966):
\begin{equation}\label{eq:mdef_gas}
  \boxed{m(p) = 2\int_{p_{\mathrm{ref}}}^{p}
         \frac{p'\,dp'}{\mu_g(p')\,Z(p')}.}
\end{equation}
Дифференциал:
\begin{equation}\label{eq:dm_gas}
  dm = \frac{2p}{\mu_g Z}\,dp,
  \qquad \nabla m = \frac{2p}{\mu_g Z}\,\nabla p.
\end{equation}

\paragraph{Пространственный оператор.}
Правая часть~\eqref{eq:gas_pde}:
\begin{equation}\label{eq:rhs_gas}
  \mathrm{div}\left(\frac{kp}{\mu_g Z}\,\nabla p\right)
  = \mathrm{div}\left(\frac{k}{2}\,\nabla m\right)
  = \frac{k}{2}\,\nabla^2 m,
\end{equation}
где последнее равенство --- при постоянной $k$. Нелинейность
\emph{полностью} поглощена заменой переменной.

\paragraph{Временной оператор.}
Левая часть~\eqref{eq:gas_pde}:
\begin{equation}\label{eq:lhs_gas}
  \frac{\partial}{\partial t}\left(\frac{p}{Z}\right)
  = \frac{1}{Z}\frac{\partial p}{\partial t}
  - \frac{p}{Z^2}\frac{\partial Z}{\partial p}\frac{\partial p}{\partial t}
  = \frac{p}{Z}\left(\frac{1}{p} - \frac{1}{Z}\frac{dZ}{dp}\right)
    \frac{\partial p}{\partial t}.
\end{equation}
Сжимаемость газа:
\begin{equation}\label{eq:cg}
  c_g = \frac{1}{\rho_g}\frac{d\rho_g}{dp}
  = \frac{Z}{p}\,\frac{d}{dp}\!\left(\frac{p}{Z}\right)
  = \frac{1}{p} - \frac{1}{Z}\frac{dZ}{dp}.
\end{equation}
Отсюда:
\begin{equation}\label{eq:lhs_gas2}
  \frac{\partial}{\partial t}\left(\frac{p}{Z}\right)
  = \frac{p\,c_g}{Z}\,\frac{\partial p}{\partial t}
  = \frac{c_g}{2}\,\frac{\mu_g Z}{1}\cdot
    \frac{2p}{\mu_g Z}\,\frac{\partial p}{\partial t}
  = \frac{\mu_g\,c_g}{2}\,\frac{\partial m}{\partial t}.
\end{equation}
Собирая:
\begin{equation}\label{eq:gas_quasi}
  \boxed{\phi\,\mu_g(p)\,c_g(p)\,\frac{\partial m}{\partial t}
  = k\,\nabla^2 m.}
\end{equation}
Пространственный оператор линеен, но коэффициент при $\partial m/\partial t$
зависит от давления через $\mu_g(p)\,c_g(p)$. Уравнение
\emph{квази}линейное.

\subsection{Замена переменной: псевдовремя}\label{ssec:pseudo_t}

Для устранения остаточной нелинейности Agarwal~(1979) ввёл псевдовремя:
\begin{equation}\label{eq:ta_def}
  \boxed{t_a = \int_0^t
    \frac{\mu_i\,c_{gi}}{\mu_g\bigl(p(\tau)\bigr)\,
          c_g\bigl(p(\tau)\bigr)}\,d\tau,}
\end{equation}
где $\mu_i$, $c_{gi}$ --- значения при начальном давлении.
Тогда $\partial m/\partial t = (\mu_i c_{gi})/(\mu_g c_g)\,
\partial m/\partial t_a$, и уравнение~\eqref{eq:gas_quasi} принимает
\emph{линейный} вид:
\begin{equation}\label{eq:gas_linear}
  \boxed{\phi\,\mu_i\,c_{gi}\,\frac{\partial m}{\partial t_a}
  = k\,\nabla^2 m.}
\end{equation}

\paragraph{Замечание.}
Строгость линеаризации --- иллюзия: в определении $t_a$ нужно знать
$p(\tau)$ в каждой точке пространства, т.\,е. решение задачи,
которую мы пытаемся получить. На практике используют осреднённое
давление $\bar{p}(t)$ из материального баланса, что корректно для
однородного пласта. Для задачи с пространственно неоднородным
профилем (как ПЗП) это приближение ухудшается.

%=============================================================================
\section{Многофазный изотермический случай (Краснов)}\label{sec:multi}
%=============================================================================

Переходим к системе Black~Oil: нефть (индекс~$o$) + свободный газ
(индекс~$g$) + вода (индекс~$w$). Капиллярное давление
пренебрежимо мало: $p_o = p_g = p_w \equiv p$.

\subsection{Уравнения неразрывности}\label{ssec:multi_cont}

\paragraph{Нефтяной компонент.}
Нефтяной компонент присутствует в двух фазах: в нефти (основная
масса) и в газе (нефтесодержание $R_v$). Масса нефтяного компонента
в единице объёма пористой среды:
\begin{equation}\label{eq:oil_mass}
  \phi\left(\rho_o(p)\,S_o + \rho_{osc}\,R_v\,\frac{S_g}{B_g}\right).
\end{equation}
Здесь $R_v(p)$ --- нефтесодержание газа (объём нефти при стандартных
условиях на объём газа при стандартных условиях), $\rho_{osc}$ ---
плотность нефти при стандартных условиях.

Поток нефтяного компонента через границу элементарного объёма:
\begin{equation}\label{eq:oil_flux}
  \rho_o\,\vec{u}_o + \rho_{osc}\,R_v\,\frac{\vec{u}_g}{B_g}.
\end{equation}
Первый член --- поток нефти в жидкой фазе. Второй --- поток нефти,
растворённой в газе: скорость газа $\vec{u}_g$ переводим в
поверхностные условия делением на $B_g$, умножаем на $R_v$ и на
$\rho_{osc}$.

Уравнение неразрывности (Краснов, ур.~0.9):
\begin{equation}\label{eq:oil_cont}
  \frac{\partial}{\partial t}\left(\phi\rho_o S_o
  + \rho_{osc}R_v\frac{\phi S_g}{B_g}\right)
  + \mathrm{div}\left(\rho_o\vec{u}_o
  + \rho_{osc}R_v\frac{\vec{u}_g}{B_g}\right) = 0.
\end{equation}

Используя $\rho(p) = \rho_{sc}/B(p)$ для каждой фазы, получим
(Краснов, ур.~0.11):
\begin{equation}\label{eq:oil_cont2}
  \frac{\partial}{\partial t}\left(\frac{\phi S_o}{B_o}
  + R_v\,\frac{\phi S_g}{B_g}\right)
  + \mathrm{div}\left(\frac{\vec{u}_o}{B_o}
  + R_v\,\frac{\vec{u}_g}{B_g}\right) = 0.
\end{equation}

\paragraph{Закон Дарси.}
\begin{equation}\label{eq:darcy}
  \vec{u}_o = -\frac{k\,k_{ro}(S_o)}{\mu_o(p)}\,\nabla p,
  \qquad
  \vec{u}_g = -\frac{k\,k_{rg}(S_g)}{\mu_g(p)}\,\nabla p.
\end{equation}

Подставляя в~\eqref{eq:oil_cont2}:
\begin{equation}\label{eq:oil_pde}
  \frac{\partial}{\partial t}\left(\frac{\phi S_o}{B_o}
  + R_v\,\frac{\phi S_g}{B_g}\right)
  - \mathrm{div}\left[k\left(\frac{k_{ro}}{B_o\mu_o}
  + R_v\,\frac{k_{rg}}{B_g\mu_g}\right)\nabla p\right] = 0.
\end{equation}
Это уравнение~(0.12) Краснова.

\subsection{Многофазное псевдодавление}\label{ssec:multi_mp}

По аналогии с однофазным случаем, введём функцию псевдодавления
для нефтяного компонента (Краснов, ур.~0.13):
\begin{equation}\label{eq:mo_def}
  \boxed{m_o(p) = \mu_{oi}\,B_{oi}\int_{p_{\mathrm{ref}}}^{p}
  \left(\frac{k_{ro}(S_o)}{B_o(p')\,\mu_o(p')}
  + R_v(p')\,\frac{k_{rg}(S_g)}{B_g(p')\,\mu_g(p')}\right)dp'.}
\end{equation}

Нормировочный множитель $\mu_{oi}B_{oi}$ обеспечивает размерность
давления и $m_o(p_i) \approx p_i$ при начальных условиях.

\paragraph{Отличие от однофазного случая.}
Подынтегральное выражение зависит не только от давления,
но и от насыщенностей $S_o$, $S_g = 1-S_o-S_w$.
Для вычисления интеграла нужна связь $S_o(p)$ ---
\textbf{замыкающее соотношение}, которого нет в общем случае.

Дифференциал:
\begin{equation}\label{eq:dmo}
  \nabla m_o = \mu_{oi}B_{oi}
  \left(\frac{k_{ro}}{B_o\mu_o}
  + R_v\frac{k_{rg}}{B_g\mu_g}\right)\nabla p.
\end{equation}

Пространственный оператор в~\eqref{eq:oil_pde}:
\begin{equation}\label{eq:rhs_multi}
  \mathrm{div}\left[k\left(\frac{k_{ro}}{B_o\mu_o}
  + R_v\frac{k_{rg}}{B_g\mu_g}\right)\nabla p\right]
  = \mathrm{div}\left(\frac{k}{\mu_{oi}B_{oi}}\nabla m_o\right)
  = \frac{1}{\mu_{oi}B_{oi}}\,\mathrm{div}(k\nabla m_o),
\end{equation}
при постоянной $k$ (или $k=k(x,y,z)$, но не $k(p)$).

\subsection{Временной оператор и общая сжимаемость}\label{ssec:multi_ct}

Раскроем временную производную в~\eqref{eq:oil_pde}. Обозначим:
\begin{equation}\label{eq:A_def}
  A(S_o, S_g, p) = \frac{\partial}{\partial p}\left[\phi\left(
  \frac{S_o}{B_o} + R_v\frac{S_g}{B_g}\right)\right]
  + \frac{\partial}{\partial S}\left[\phi\left(
  \frac{S_o}{B_o} + R_v\frac{S_g}{B_g}\right)\right]
  \frac{dS}{dp},
\end{equation}
так что:
\begin{equation}\label{eq:time_A}
  \frac{\partial}{\partial t}\left(\frac{\phi S_o}{B_o}
  + R_v\frac{\phi S_g}{B_g}\right)
  = A(S_o, S_g, p)\,\frac{\partial p}{\partial t}.
\end{equation}
Это уравнение~(0.19)--(0.20) Краснова.

\paragraph{Подробное раскрытие~$A$.}
Используем определения сжимаемостей фаз (Краснов, ур.~0.18):
\begin{equation}\label{eq:compress}
  C_o = -\frac{1}{B_o}\frac{dB_o}{dp} = -\frac{B_o'}{B_o},
  \qquad
  C_g = -\frac{B_g'}{B_g},
  \qquad
  C_\phi = \frac{1}{\phi}\frac{d\phi}{dp}.
\end{equation}
Раскрывая~\eqref{eq:A_def} по слагаемым (здесь опускаем
сжимаемость породы $C_\phi$ --- она мала по сравнению с
сжимаемостью флюида в условиях разгазирования):

\begin{align}
  \frac{\partial}{\partial p}\left(\frac{S_o}{B_o}\right)
  &= -\frac{S_o B_o'}{B_o^2} = \frac{S_o C_o}{B_o},
  \label{eq:dSoBo_dp} \\
  \frac{\partial}{\partial p}\left(R_v\frac{S_g}{B_g}\right)
  &= R_v'\frac{S_g}{B_g} + R_v\frac{S_g C_g}{B_g},
  \label{eq:dRvSgBg_dp} \\
  \frac{\partial}{\partial S_o}\left(\frac{S_o}{B_o}
  + R_v\frac{S_g}{B_g}\right)
  &= \frac{1}{B_o} - \frac{R_v}{B_g}
  \qquad (\text{т.\,к. } S_g = 1 - S_o - S_w).
  \label{eq:dSoBo_dS}
\end{align}

Собирая:
\begin{equation}\label{eq:A_full}
  A = \phi\left[\frac{S_o C_o}{B_o}
  + R_v'\frac{S_g}{B_g} + R_v\frac{S_g C_g}{B_g}
  + \left(\frac{1}{B_o} - \frac{R_v}{B_g}\right)\frac{dS_o}{dp}
  \right].
\end{equation}

\paragraph{Общая сжимаемость двухфазной системы.}
На малых временах (или при медленном изменении насыщенности)
связь $dS_o/dp$ определяется из баланса (Краснов, ур.~0.77):
\begin{equation}\label{eq:dSdp_small}
  \frac{dS_o}{dp} = S_o\frac{B_o'}{B_o}
  + \frac{\lambda_o}{\lambda_t}\,C_t,
\end{equation}
где
\begin{equation}\label{eq:Ct}
  \boxed{C_t = S_o\,C_o + S_g\,C_g + S_o\,R_s'\,\frac{B_g}{B_o}}
\end{equation}
--- общая сжимаемость системы (Краснов, ур.~0.80--0.81),
\begin{equation}\label{eq:lambda}
  \lambda_o = \frac{k_{ro}}{\mu_o},
  \qquad
  \lambda_g = \frac{k_{rg}}{\mu_g},
  \qquad
  \lambda_t = \lambda_o + \lambda_g.
\end{equation}

Член $C_{R_s} = S_o R_s' B_g/B_o$ в~\eqref{eq:Ct} отражает
\emph{фазовый переход}: при снижении давления газ выделяется из
нефти ($R_s' > 0$), что эквивалентно аномально высокой сжимаемости.
Отношение $B_g/B_o$ учитывает разницу плотностей газа в свободном
и растворённом состояниях.

\subsection{Линеаризованное уравнение}\label{ssec:multi_lin}

Подставляя~\eqref{eq:time_A} и~\eqref{eq:rhs_multi}
в~\eqref{eq:oil_pde}:
\begin{equation}\label{eq:oil_Am}
  \mu_{oi}B_{oi}\,A\,\frac{\partial p}{\partial t}
  = \mathrm{div}(k\nabla m_o).
\end{equation}

Из определения $m_o$:
\begin{equation}\label{eq:dmodt}
  \frac{\partial m_o}{\partial t}
  = \mu_{oi}B_{oi}\left(\frac{k_{ro}}{B_o\mu_o}
  + R_v\frac{k_{rg}}{B_g\mu_g}\right)\frac{\partial p}{\partial t}.
\end{equation}

Обозначим подвижность нефтяного компонента:
\begin{equation}\label{eq:alpha_o}
  \alpha_o(S_o, p) = \frac{k_{ro}}{B_o\mu_o}
  + R_v\frac{k_{rg}}{B_g\mu_g}.
\end{equation}

Тогда $\partial m_o/\partial t = \mu_{oi}B_{oi}\,\alpha_o\,
\partial p/\partial t$, и из~\eqref{eq:oil_Am}:
\begin{equation}\label{eq:oil_nonlin}
  \frac{A}{\alpha_o}\,\frac{\partial m_o}{\partial t}
  = \mathrm{div}(k\nabla m_o).
\end{equation}

Это уравнение~(0.21) Краснова. Отношение $A/\alpha_o$ ---
обобщённый коэффициент сжимаемости-к-подвижности двухфазной системы,
аналог $\phi\mu_g c_g/k$ в однофазном случае.

\subsection{Многофазное псевдовремя}\label{ssec:multi_pt}

По аналогии с Agarwal, вводим псевдовремя (Краснов, ур.~0.22):
\begin{equation}\label{eq:tao_def}
  \boxed{t_{ao} = \mu_{oi}\,C_{ti}\,\phi\int_0^t
  \frac{\alpha_o(\tau)}{A(\tau)}\,d\tau,}
\end{equation}
где $C_{ti}$ --- начальная общая сжимаемость. Тогда:
\begin{equation}\label{eq:oil_linear}
  \boxed{\mu_{oi}\,C_{ti}\,\phi\,
  \frac{\partial m_o}{\partial t_{ao}}
  = \mathrm{div}(k\nabla m_o).}
\end{equation}
Уравнение \emph{формально} линейное --- аналог~\eqref{eq:gas_linear}.

\paragraph{Оговорка.}
Линейность формальна по тем же причинам, что и в однофазном случае:
$t_{ao}$ зависит от решения (через $S_o(t)$ и $p(t)$). Кроме того,
$m_o$ определено через неизвестную связь $S_o(p)$. Подход работает
при наличии замыкающего соотношения (GOR~=~const, малые времена
и т.\,д.).


%=============================================================================
\section{Неизотермическое обобщение}\label{sec:noniso}
%=============================================================================

\subsection{Постановка вопроса}\label{ssec:question}

В задаче ПЗП температура --- неизвестная, определяемая балансом
энергии. PVT-свойства зависят от двух переменных: $p$ и $T$.
В частности, $R_s = R_s(p,T)$, $B_o = B_o(p,T)$, $\mu_o = \mu_o(p,T)$.

Вопрос: можно ли ввести функции $m_o(p,T)$ и $t_{ao}(t)$ так, чтобы
система уравнений (фильтрация + энергия) допускала хотя бы частичную
линеаризацию?

\subsection{Уравнение фильтрации при $T \neq \mathrm{const}$}
\label{ssec:filt_T}

Уравнение~\eqref{eq:oil_pde} сохраняет структуру, но все
PVT-свойства теперь зависят от $(p,T)$:
\begin{equation}\label{eq:oil_pde_T}
  \frac{\partial}{\partial t}\left(\frac{\phi S_o}{B_o(p,T)}
  + R_v(p,T)\frac{\phi S_g}{B_g(p,T)}\right)
  - \mathrm{div}\left[k\left(\frac{k_{ro}}{B_o\mu_o}
  + R_v\frac{k_{rg}}{B_g\mu_g}\right)\nabla p\right] = 0.
\end{equation}

Попытка ввести псевдодавление:
\begin{equation}\label{eq:mo_T}
  m_o(p,T) = \mu_{oi}B_{oi}\int_{p_{\mathrm{ref}}}^{p}
  \left(\frac{k_{ro}}{B_o(p',T)\mu_o(p',T)}
  + R_v(p',T)\frac{k_{rg}}{B_g(p',T)\mu_g(p',T)}\right)dp'.
\end{equation}

\paragraph{Проблема 1: $m_o$ --- функция двух переменных.}
В изотермическом случае $m_o = m_o(p)$ --- скалярная функция одной
переменной, и $\nabla m_o = m_o'(p)\nabla p$. В неизотермическом
случае:
\begin{equation}\label{eq:grad_mo_T}
  \nabla m_o = \frac{\partial m_o}{\partial p}\,\nabla p
  + \frac{\partial m_o}{\partial T}\,\nabla T.
\end{equation}
Пространственный оператор $\mathrm{div}(k\nabla m_o)$ содержит
$\nabla T$, что \emph{связывает} уравнение фильтрации с уравнением
энергии уже на уровне пространственного оператора.

\paragraph{Проблема 2: верхний предел интеграла.}
В~\eqref{eq:mo_T} интегрирование ведётся по $p'$ при фиксированном
$T$. Но температура в пласте неоднородна: в разных точках
интеграл вычисляется при разных $T$. Таким образом, $m_o$ ---
нелокальная характеристика: её значение в точке $(r,t)$ зависит
от температуры \emph{в этой же точке}.

\subsection{Уравнение энергии}\label{ssec:energy}

Уравнение баланса энергии многофазной системы в пористой среде
(см.\ Draft4/two\_phase\_energy.tex):
\begin{multline}\label{eq:energy}
  \frac{\partial}{\partial t}\bigl[
  \phi(\rho_o c_{p,o} S_o + \rho_g c_{p,g} S_g)T
  + (1-\phi)(\rho c_p)_s T\bigr] \\
  + \mathrm{div}\bigl[(\rho_o c_{p,o}\vec{u}_o
  + \rho_g c_{p,g}\vec{u}_g)\,T\bigr]
  = \mathrm{div}(\lambda_{\mathrm{eff}}\nabla T)
  + Q_{\mathrm{JT}} + Q_{\mathrm{deg}},
\end{multline}
где $Q_{\mathrm{JT}}$ --- нагрев/охлаждение Джоуля--Томсона,
$Q_{\mathrm{deg}}$ --- скрытая теплота разгазирования.

Это уравнение не содержит $m_o$ и не упрощается введением
псевдодавления.

\subsection{Анализ возможности линеаризации}\label{ssec:analysis}

Рассмотрим, что мешает обобщению:

\begin{enumerate}[label=(\alph*)]
  \item \textbf{Смешанный градиент.}
        Из~\eqref{eq:grad_mo_T}: $\mathrm{div}(k\nabla m_o)$
        содержит и $\nabla p$, и $\nabla T$. Замена $p \to m_o$
        не развязывает систему --- она вносит $\nabla T$
        в уравнение фильтрации.

  \item \textbf{Температура в подынтегральном выражении.}
        Даже если рассматривать $m_o$ при \emph{локальной}
        температуре, функция $m_o(p,T(r,t))$ зависит от
        пространственно-временного поля $T(r,t)$.

  \item \textbf{Уравнение энергии остаётся нелинейным.}
        Введение $m_o$ не упрощает~\eqref{eq:energy}:
        источники $Q_{\mathrm{JT}}$ и $Q_{\mathrm{deg}}$
        зависят от $\partial p/\partial t$, $\nabla p$,
        $\partial S/\partial t$ --- все эти величины
        не выражаются через $m_o$ без дополнительных допущений.

  \item \textbf{Насыщенность зависит от двух переменных.}
        В изотермическом случае $S = S(p)$ при фиксированном GOR.
        При $T \neq \mathrm{const}$: $R_s = R_s(p,T)$, и
        $S = S(p,T)$ --- замыкание через GOR~=~const даёт
        неявную зависимость $S(p,T)$ вместо $S(p)$.
\end{enumerate}

\subsection{Предварительный вывод}\label{ssec:conclusion}

Полная линеаризация системы фильтрация~+ энергия через
псевдодавление \textbf{невозможна}: замена переменной $p \to m_o$
не развязывает уравнения и вносит дополнительную нелинейность через
$\nabla T$.

Однако возможны \emph{частичные} упрощения:

\begin{enumerate}[nosep]
  \item \textbf{Слабая неизотермичность.}
        Если $\Delta T / T_0 \ll 1$ (типично $\Delta T \sim 1$--$5$~К
        при $T_0 \sim 350$~К), то зависимость PVT-свойств от $T$
        можно линеаризовать:
        $B_o(p,T) \approx B_o(p,T_0)(1 + \alpha_B\,\delta T)$.
        В этом приближении $\partial m_o/\partial T \sim O(\delta T)$
        и член $(\partial m_o/\partial T)\nabla T$ в~\eqref{eq:grad_mo_T}
        является поправкой второго порядка.

  \item \textbf{Разделение масштабов.}
        Если тепловой фронт много шире гидродинамического
        ($\lambda_{\mathrm{eff}} \gg \phi\rho c_p u$),
        температура --- медленная переменная. Тогда на каждом
        временном шаге можно считать $T(r)$ известной и строить
        $m_o(p; T(r))$ как в изотермическом случае, но
        с <<замороженной>> температурой.

  \item \textbf{Псевдодавление при фиксированной температуре.}
        Определить $m_o(p) = m_o(p; T(r,t))$ при локальной
        температуре. Тогда пространственный оператор линеаризуется
        с точностью до $O(\nabla T)$-поправки, а уравнение энергии
        решается отдельно.
\end{enumerate}

Эти варианты требуют оценки масштабов --- какие члены велики
в условиях ПЗП. Это предмет следующего раздела.


%=============================================================================
\section{Оценка масштабов для условий ПЗП}\label{sec:scales}
%=============================================================================

\textit{(Раздел будет дополнен.)}


%=============================================================================
\section*{Список литературы}
%=============================================================================

\begin{enumerate}[label={[\arabic*]}, nosep]
  \item Al-Hussainy R., Ramey H.J.\,Jr., Crawford P.B.
        The Flow of Real Gases Through Porous Media.
        \textit{J.~Pet.~Tech.}, 1966, \textbf{18}(5), 624--636.
        SPE~1243-A-PA.
  \item Agarwal R.G. <<Real Gas Pseudo-Time>> --- A New Function
        for Pressure Buildup Analysis of MHF Gas Wells.
        SPE~8279, 1979.
  \item Краснов~В.А. (предп.). Текущий (flowing) материальный
        баланс для многофазной фильтрации. Внутренний методический
        документ, б/г.
\end{enumerate}

\end{document}
